% Shared exercise chunk: l4-sharing-horizon
\begin{exercisechunk}
\item \textbf{Sequence length and the training budget.\suppmark}
\label[exercise]{ex:sharing-horizon}
Use the grouped Bernoulli model and smoothed estimator from
\cref{sec:parameter-sharing}. Assume \(p_j\in[a,1-a]\), with
\(0<a<1/2\).
\begin{enumerate}[(a),beginpenalty=10000]
\item Keep \(n\) fixed and suppose \(p_t\in[a,1-a]\).
Use \cref{ex:sharing}(b) to give explicit upper and lower bounds proportional to
\(T/n\) for the unshared model's expected sequence KL.
Compare this growth with the shared model's limit in
\cref{eq:shared-bernoulli}.
\item Use one new parameter for each successive block of lengths
\(1,2,3,\ldots\). Compute \(d(T)\) at complete-block endpoints and
interpret \cref{eq:sharing-rate,eq:sharing-leading}. How can normalized
KL improve while sequence KL grows?
\item Fix the total training-token budget \(L=nT\) and use
\(\kappa'_L(T)=\kappa_{L/T}(T)\) to repeat the shared and unshared
leading-order comparisons. Explain why increasing \(T\) worsens
expected sequence KL in the shared model. Compare the numbers of
observations per parameter in the two models. Translate the expected-KL bounds into
bounds on the discrepancy of a cost in \([0,1]\).
\end{enumerate}

\end{exercisechunk}
